% GENERATED FILE. Edit canonical lessons, then run npm run generate:books.
% canonical-source-hash: fnv1a64:11d6c7455cf2e3cc
% canonical-lessons: GE-C318-krieg, GE-C318-frieden, GE-C318-schlacht, GE-C318-besiegen, GE-C318-niederlage

\chapter{Krieg und Frieden}
\label{ch:ge-krieg-und-frieden}
\hlchaptermodality{318}

\begin{chapteropening}
\textbf{By the end of this chapter:} \emph{I can tell about history: war, peace, a battle, to defeat and a defeat.}

Complete the last lesson of chapter 318: I can tell about history: war, peace, a battle, to defeat and a defeat.
\end{chapteropening}

\section[der Krieg]{der Krieg — the war}

\label{lesson:GE-C318-krieg}



\begin{quote}
Before the new one: say the German for to hand down, then the German for the fate.
\end{quote}

\subsection*{You'll want to know: der Krieg}
\textbf{der Krieg} — \textquotedblleft{}the war\textquotedblright{}.

\begin{quote}
Der Krieg dauerte sechs Jahre.
\end{quote}

\begin{quote}
\emph{The war lasted six years.}
\end{quote}

Plural: \emph{die Kriege}.

\begin{cousinweb}
\emph{Krieg} is inherited from Middle High German \emph{kriec}, Old High German \emph{krieg}, from Proto-Germanic \emph{krīgaz} (EtymDB); it is usually said to have first meant \textquotedblleft{}effort, struggle\textquotedblright{}.
\end{cousinweb}

\begin{grammarlens}[title={What you've built}]
One more word for telling, explaining and describing at B1.
\end{grammarlens}

\subsection*{Your turn}
\begin{itemize}
\raggedright
  \item \emph{Say it:} \emph{der Krieg}
  \item \emph{Say it:} \emph{der Krieg}, once more
  \item \emph{Say it:} a sentence of your own with \emph{der Krieg}
  \item \emph{Recall:} say the German for to hand down, then the German for the fate, then say \emph{der Krieg} again
\end{itemize}

\begin{tcolorbox}[breakable,colback=teal!4,colframe=teal!35!black,title={Before you move on}]
What does der Krieg mean? (\textquotedblleft{}The war\textquotedblright{}.) Say it once more.
\end{tcolorbox}
\section[der Frieden]{der Frieden — the peace}

\label{lesson:GE-C318-frieden}



\begin{quote}
Before the new one: say the German for the fate, then the German for the war.
\end{quote}

\subsection*{You'll want to know: der Frieden}
\textbf{der Frieden} — \textquotedblleft{}the peace\textquotedblright{}.

\begin{quote}
Nach dem Krieg kam endlich Frieden.
\end{quote}

\begin{quote}
\emph{After the war peace came at last.}
\end{quote}

Also \emph{der Friede}, older and more formal.

\begin{cousinweb}
\emph{Frieden} goes back to Proto-Germanic \emph{friþuz} \textquotedblleft{}peace\textquotedblright{} (EtymDB), with the old English cousin \emph{frith} \textquotedblleft{}peace\textquotedblright{}.
\end{cousinweb}

\begin{grammarlens}[title={What you've built}]
One more word for telling, explaining and describing at B1.
\end{grammarlens}

\subsection*{Your turn}
\begin{itemize}
\raggedright
  \item \emph{Say it:} \emph{der Frieden}
  \item \emph{Say it:} \emph{der Frieden}, once more
  \item \emph{Say it:} a sentence of your own with \emph{der Frieden}
  \item \emph{Recall:} say the German for the fate, then the German for the war, then say \emph{der Frieden} again
\end{itemize}

\begin{tcolorbox}[breakable,colback=teal!4,colframe=teal!35!black,title={Before you move on}]
What does der Frieden mean? (\textquotedblleft{}The peace\textquotedblright{}.) Say it once more.
\end{tcolorbox}
\section[die Schlacht]{die Schlacht — the battle}

\label{lesson:GE-C318-schlacht}



\begin{quote}
Before the new one: say the German for the war, then the German for the peace.
\end{quote}

\subsection*{You'll want to know: die Schlacht}
\textbf{die Schlacht} — \textquotedblleft{}the battle\textquotedblright{}.

\begin{quote}
Die Schlacht fand vor der Stadt statt.
\end{quote}

\begin{quote}
\emph{The battle took place outside the town.}
\end{quote}

Plural: \emph{die Schlachten}.

\begin{cousinweb}
\emph{Schlacht} is inherited from Old High German \emph{slahta}, Proto-Germanic \emph{slahtō} \textquotedblleft{}a striking\textquotedblright{} (EtymDB), from the root of \emph{schlagen} \textquotedblleft{}to strike\textquotedblright{}.
\end{cousinweb}

\begin{grammarlens}[title={What you've built}]
One more word for telling, explaining and describing at B1.
\end{grammarlens}

\subsection*{Your turn}
\begin{itemize}
\raggedright
  \item \emph{Say it:} \emph{die Schlacht}
  \item \emph{Say it:} \emph{die Schlacht}, once more
  \item \emph{Say it:} a sentence of your own with \emph{die Schlacht}
  \item \emph{Recall:} say the German for the war, then the German for the peace, then say \emph{die Schlacht} again
\end{itemize}

\begin{tcolorbox}[breakable,colback=teal!4,colframe=teal!35!black,title={Before you move on}]
What does die Schlacht mean? (\textquotedblleft{}The battle\textquotedblright{}.) Say it once more.
\end{tcolorbox}
\section[besiegen]{besiegen — to defeat, to beat}

\label{lesson:GE-C318-besiegen}



\begin{quote}
Before the new one: say the German for the peace, then the German for the battle.
\end{quote}

\subsection*{You'll want to know: besiegen}
\textbf{besiegen} — \textquotedblleft{}to defeat, to beat\textquotedblright{}.

\begin{quote}
Die kleine Armee besiegte die große.
\end{quote}

\begin{quote}
\emph{The small army defeated the big one.}
\end{quote}

Regular: \emph{besiegte}, \emph{hat besiegt}.

\begin{cousinweb}
\emph{besiegen} is built on \emph{Sieg} \textquotedblleft{}victory\textquotedblright{}, Old High German \emph{sigu}, Proto-Germanic \emph{segaz} (EtymDB), with \emph{be-}.
\end{cousinweb}

\begin{grammarlens}[title={What you've built}]
One more word for telling, explaining and describing at B1.
\end{grammarlens}

\subsection*{Your turn}
\begin{itemize}
\raggedright
  \item \emph{Say it:} \emph{besiegen}
  \item \emph{Say it:} \emph{besiegen}, once more
  \item \emph{Say it:} a sentence of your own with \emph{besiegen}
  \item \emph{Recall:} say the German for the peace, then the German for the battle, then say \emph{besiegen} again
\end{itemize}

\begin{tcolorbox}[breakable,colback=teal!4,colframe=teal!35!black,title={Before you move on}]
What does besiegen mean? (\textquotedblleft{}To defeat, to beat\textquotedblright{}.) Say it once more.
\end{tcolorbox}
\section[die Niederlage]{die Niederlage — the defeat}

\label{lesson:GE-C318-niederlage}



\begin{quote}
Before the new one: say the German for the battle, then the German for to defeat.
\end{quote}

\subsection*{You'll want to know: die Niederlage}
\textbf{die Niederlage} — \textquotedblleft{}the defeat\textquotedblright{}.

\begin{quote}
Nach der Niederlage gingen die Soldaten nach Hause.
\end{quote}

\begin{quote}
\emph{After the defeat the soldiers went home.}
\end{quote}

Plural: \emph{die Niederlagen}. Also in sport.

\begin{cousinweb}
\emph{Niederlage} comes from Middle High German \emph{niderlāge} (EtymDB), from \emph{nieder} \textquotedblleft{}down\textquotedblright{}, Old High German \emph{nidar}, and \emph{Lage} \textquotedblleft{}a lying, a laying\textquotedblright{}, usually linked to \emph{legen} \textquotedblleft{}to lay\textquotedblright{}: being laid low.
\end{cousinweb}

\begin{grammarlens}[title={What you've built}]
One more word for telling, explaining and describing at B1.
\end{grammarlens}

\subsection*{Your turn}
\begin{itemize}
\raggedright
  \item \emph{Say it:} \emph{die Niederlage}
  \item \emph{Say it:} \emph{die Niederlage}, once more
  \item \emph{Say it:} a sentence of your own with \emph{die Niederlage}
  \item \emph{Recall:} say the German for the battle, then the German for to defeat, then say \emph{die Niederlage} again
\end{itemize}

\begin{tcolorbox}[breakable,colback=teal!4,colframe=teal!35!black,title={Before you move on}]
What does die Niederlage mean? (\textquotedblleft{}The defeat\textquotedblright{}.) Say it once more.
\end{tcolorbox}
